% ============================================================
%  Capítulo -- Costes
% ============================================================

\chapter{Costes}
\label{sec:costes}

Este capítulo presenta una estimación económica del proyecto OmniLitter. Al
tratarse de un Trabajo de Fin de Grado, una parte importante de los recursos no ha
supuesto un desembolso directo, ya que se han utilizado equipos personales,
servicios con planes gratuitos y herramientas disponibles para uso académico. No
obstante, resulta útil valorar el coste equivalente del proyecto para comprender
su dimensión real y justificar el esfuerzo necesario para construir una plataforma
compuesta por backend, portal web, aplicación móvil, documentación y despliegue.

Las estimaciones se han realizado diferenciando entre costes de personal,
infraestructura, herramientas software, equipamiento y otros costes operativos. En
los casos en los que no existió un gasto directo, se indica como coste imputado o
amortizado, ya que representa el valor económico aproximado de los recursos
utilizados durante el desarrollo.

\section{Costes de personal}

El coste principal del proyecto corresponde al tiempo dedicado por los dos
integrantes del equipo. Aunque no se trata de un coste salarial real, se ha
estimado su valor económico utilizando una tarifa de referencia de
\textbf{15 EUR/hora}, razonable para un perfil junior o en prácticas dentro de un
proyecto académico de ingeniería del software.

La valoración de horas se basa en la dedicación real registrada durante el
proyecto y no en la estimación inicial de planificación. Incluye análisis de
requisitos, diseño, implementación de los tres módulos principales, integración,
pruebas, despliegue, elaboración de la memoria, revisión final y preparación de
los materiales de defensa. La dedicación registrada entre los dos integrantes del
equipo asciende a \textbf{603,34 horas} de trabajo.

\begin{table}[H]
\centering
\caption{Costes de personal imputados a partir de horas reales por actividad}
\label{tab:costes_personal}
\small
\renewcommand{\arraystretch}{1.18}
\begin{tabularx}{\textwidth}{p{4.2cm}p{2.2cm}p{2.4cm}X}
\toprule
\textbf{Actividad} & \textbf{Horas reales} & \textbf{Coste imputado} & \textbf{Justificación} \\
\midrule
Planificación inicial & 42,50 h & 637,50 EUR & Incluye reuniones iniciales, análisis del alcance, organización del backlog y planificación general. \\
Sprint inicial & 64,75 h & 971,25 EUR & Arquitectura, configuración, CI/CD, repositorio, Docker, Prisma y estructura inicial del proyecto. \\
Sprint 1 & 78,25 h & 1.173,75 EUR & Base funcional del sistema: autenticación, usuarios, grupos, campañas, sesiones y primeras observaciones. \\
Sprint 2 & 84,50 h & 1.267,50 EUR & Captura de campo, fotografías, GPS, validaciones, persistencia local inicial y consulta web. \\
Sprint 3 & 91,75 h & 1.376,25 EUR & Sincronización, mapas, métricas, exportaciones y consolidación del flujo offline. \\
Sprint 4 & 72,00 h & 1.080,00 EUR & Cierre funcional, administración, ajustes finales, pruebas regresivas y estabilización. \\
Mockups y diseño UI/UX & 48,75 h & 731,25 EUR & Diseño y revisión de prototipos del portal web y de la aplicación móvil. \\
Despliegue y preparación de releases & 36,51 h & 547,65 EUR & Configuración de entornos, publicación de versiones, builds y preparación de entregables. \\
Memoria y documentación & 74,25 h & 1.113,75 EUR & Redacción de memoria, anexos, manuales, evidencias y revisión documental. \\
Revisión final, vídeo y presentación & 10,08 h & 151,20 EUR & Correcciones finales de la memoria, preparación del vídeo demostrativo y elaboración de la presentación de defensa. \\
\midrule
\textbf{Total} & \textbf{603,34 h} & \textbf{9.050,10 EUR} & Coste imputado del esfuerzo de desarrollo del equipo. \\
\bottomrule
\end{tabularx}
\end{table}

Esta partida concentra la mayor parte del presupuesto porque OmniLitter es un
proyecto intensivo en desarrollo software. La existencia de varios módulos
conectados entre sí incrementó el esfuerzo de integración, especialmente en la
sincronización \textit{offline-first} y en la coordinación entre backend, web y
móvil.

\section{Costes de infraestructura y despliegue}

La infraestructura utilizada durante el desarrollo se apoyó principalmente en
servicios con planes gratuitos o recursos de bajo coste. GitHub se empleó para el
repositorio, gestión de issues, tablero Kanban, releases y automatizaciones.
Además, el despliegue se planteó con servicios adecuados para un proyecto
académico, como Vercel, Render, Railway o entornos equivalentes.

Durante el desarrollo, estos servicios no generaron un coste directo relevante. A
efectos de estimación, se contempla un escenario de uso académico durante seis
meses, con infraestructura en planes gratuitos y sin contratación de dominio
propio.

\begin{table}[H]
\centering
\caption{Estimación de costes de infraestructura}
\label{tab:costes_infraestructura}
\small
\renewcommand{\arraystretch}{1.18}
\begin{tabularx}{\textwidth}{p{4.0cm}p{2.5cm}p{2.5cm}X}
\toprule
\textbf{Recurso} & \textbf{Periodo} & \textbf{Coste estimado} & \textbf{Justificación} \\
\midrule
Repositorio GitHub y GitHub Projects & 6 meses & 0 EUR & Uso de funcionalidades gratuitas para repositorio, issues, tablero Kanban, Pull Requests y releases. \\
GitHub Actions & 6 meses & 0 EUR & Uso dentro de los límites gratuitos disponibles para automatizaciones ligeras de CI/CD. \\
Hosting portal web & 6 meses & 0 EUR & Despliegue compatible con planes gratuitos de servicios como Vercel o equivalentes. \\
Backend y API & 6 meses & 0 EUR & Entorno de despliegue académico o gratuito, suficiente para validación y demostración. \\
Base de datos PostgreSQL & 6 meses & 0 EUR & Uso de instancia gratuita o entorno local durante desarrollo y pruebas. \\
Dominio propio & No contratado & 0 EUR & No fue necesario adquirir un dominio para la entrega académica del proyecto. \\
\midrule
\textbf{Total} & -- & \textbf{0 EUR} & Sin gasto directo de infraestructura durante el desarrollo académico. \\
\bottomrule
\end{tabularx}
\end{table}

En un escenario de operación real, estos costes podrían aumentar por la
contratación de una base de datos persistente, almacenamiento de fotografías,
dominio propio, copias de seguridad y mayor capacidad de cómputo. Sin embargo,
para el alcance del trabajo de fin de grado, los planes gratuitos fueron suficientes para validar la
solución.

\section{Costes de herramientas software}

Las herramientas de desarrollo y documentación seleccionadas se eligieron también
con criterios de coste. Se priorizaron tecnologías libres, servicios gratuitos o
licencias disponibles para estudiantes. Esto permitió reducir el presupuesto sin
limitar la calidad técnica del producto.

\begin{table}[H]
\centering
\caption{Estimación de costes de herramientas software}
\label{tab:costes_herramientas}
\small
\renewcommand{\arraystretch}{1.18}
\begin{tabularx}{\textwidth}{p{4.0cm}p{2.7cm}X}
\toprule
\textbf{Herramienta} & \textbf{Coste estimado} & \textbf{Uso en el proyecto} \\
\midrule
Visual Studio Code & 0 EUR & Entorno principal de desarrollo y edición de código. \\
Git y GitHub & 0 EUR & Control de versiones, repositorio centralizado, issues, Pull Requests y releases. \\
Node.js, TypeScript, React, Express y Expo & 0 EUR & Stack principal de desarrollo, basado en tecnologías de código abierto. \\
PostgreSQL y SQLite & 0 EUR & Persistencia remota y local mediante motores gratuitos. \\
Overleaf & 0 EUR & Redacción inicial de la memoria mediante plan gratuito. \\
LaTeX local & 0 EUR & Compilación final de la memoria tras las limitaciones de Overleaf. \\
Microsoft Teams, Outlook, Word y OneDrive & 0 EUR & Herramientas disponibles mediante cuentas académicas o personales para comunicación y documentación. \\
\midrule
\textbf{Total} & \textbf{0 EUR} & No se contrataron licencias específicas para el proyecto. \\
\bottomrule
\end{tabularx}
\end{table}

La ausencia de coste de licencias fue posible porque el proyecto se apoyó en un
ecosistema ampliamente disponible para desarrollo web y móvil. La principal
limitación no fue económica, sino de aprendizaje, configuración y coordinación
entre herramientas.

\section{Costes de equipamiento}

El desarrollo se realizó con equipos personales ya disponibles. Por ello, no hubo
un gasto de adquisición específico. Aun así, para valorar económicamente el uso de
estos recursos se aplica una estimación por amortización, considerando el tiempo de
uso durante el proyecto frente a la vida útil aproximada de los dispositivos.

Se asume un periodo de proyecto de seis meses, una vida útil de cuatro años para
los ordenadores portátiles y tres años para el dispositivo móvil utilizado en
pruebas. Esta estimación no representa dinero desembolsado, sino el coste
proporcional de uso del equipamiento.

\begin{table}[H]
\centering
\caption{Estimación de costes de equipamiento amortizado}
\label{tab:costes_equipamiento}
\small
\renewcommand{\arraystretch}{1.18}
\begin{tabularx}{\textwidth}{p{3.7cm}p{2.0cm}p{2.7cm}p{2.6cm}X}
\toprule
\textbf{Recurso} & \textbf{Unidades} & \textbf{Valor unitario} & \textbf{Coste imputado} & \textbf{Justificación} \\
\midrule
Ordenadores portátiles & 2 & 900 EUR & 225 EUR & Amortización de 6 meses sobre una vida útil estimada de 48 meses. \\
Dispositivo Android de pruebas & 1 & 300 EUR & 50 EUR & Amortización de 6 meses sobre una vida útil estimada de 36 meses. \\
Conexión a Internet y electricidad & 1 & -- & 60 EUR & Estimación proporcional del consumo y conectividad durante el trabajo. \\
\midrule
\textbf{Total} & -- & -- & \textbf{335 EUR} & Coste imputado de recursos físicos y operativos. \\
\bottomrule
\end{tabularx}
\end{table}

Estos costes son reducidos en comparación con el esfuerzo humano, ya que el
proyecto no requirió hardware específico ni servidores propios. El elemento más
relevante fue disponer de al menos un dispositivo móvil real para validar
funcionalidades como GPS, cámara, almacenamiento local y sincronización.

\section{Plan de contingencia}

Además de los costes estimados, se contempla una reserva de contingencia para
cubrir desviaciones razonables durante el desarrollo. En un proyecto de estas
características, los principales factores de incertidumbre no proceden de la
compra de licencias o infraestructura, sino del tiempo adicional necesario para
resolver problemas técnicos, adaptar requisitos o estabilizar la integración entre
módulos.

La contingencia se calcula como un \textbf{10\%} sobre el coste base estimado del
proyecto, excluyendo únicamente costes no previstos que requerirían una decisión
externa, como contratar un servicio de pago o adquirir nuevo equipamiento. Esta
reserva se considera adecuada para un trabajo académico con alcance controlado,
pero con riesgos técnicos relevantes asociados al funcionamiento offline, el
despliegue y la coordinación entre backend, web y aplicación móvil.

\begin{table}[H]
\centering
\caption{Plan de contingencia del proyecto}
\label{tab:plan_contingencia}
\small
\renewcommand{\arraystretch}{1.18}
\begin{tabularx}{\textwidth}{p{4.1cm}p{1.5cm}p{2.2cm}X}
\toprule
\textbf{Riesgo cubierto} & \textbf{Prob.} & \textbf{Reserva} & \textbf{Aplicación prevista} \\
\midrule
Cambios de requisitos del cliente & Media & 295 EUR & Revisión de funcionalidades, ajustes de alcance, adaptación del backlog y pequeñas modificaciones de flujo. \\
Complejidad de sincronización offline & Alta & 345 EUR & Tiempo adicional para depurar estados locales, reintentos, conflictos y validación de datos sincronizados. \\
Problemas de integración entre módulos & Media & 148 EUR & Corrección de incompatibilidades entre backend, portal web, aplicación móvil y contratos compartidos. \\
Incidencias de despliegue o configuración & Media & 85 EUR & Ajustes de variables de entorno, builds, servicios externos y documentación operativa. \\
Margen general de cierre & Baja & 65,51 EUR & Correcciones menores, revisión final, pruebas regresivas y ajustes documentales de última hora. \\
\midrule
\textbf{Total contingencia} & -- & \textbf{938,51 EUR} & Reserva equivalente al 10\% del coste base del proyecto. \\
\bottomrule
\end{tabularx}
\end{table}

Este plan no implica que la contingencia deba consumirse obligatoriamente. Su
función es reflejar una práctica de planificación prudente: reservar margen ante
riesgos identificados, evitar que una desviación técnica afecte al cierre del
proyecto y disponer de una valoración económica más realista si OmniLitter se
plantease como un desarrollo profesional.

\section{Costes globales del proyecto}

El coste global se obtiene sumando las partidas anteriores y la reserva de
contingencia definida. La mayor parte del importe corresponde al esfuerzo humano,
mientras que infraestructura y licencias presentan coste directo nulo gracias al
uso de herramientas gratuitas o ya disponibles.

\begin{table}[H]
\centering
\caption{Resumen global de costes del proyecto}
\label{tab:costes_globales}
\small
\renewcommand{\arraystretch}{1.18}
\begin{tabularx}{0.85\textwidth}{X>{\raggedleft\arraybackslash}p{3.0cm}}
\toprule
\textbf{Categoría de coste} & \textbf{Importe estimado} \\
\midrule
Costes de personal & 9.050,10 EUR \\
Costes de infraestructura y despliegue & 0 EUR \\
Costes de herramientas software & 0 EUR \\
Costes de equipamiento y operación & 335 EUR \\
Reserva de contingencia & 938,51 EUR \\
\midrule
\textbf{Coste total estimado} & \textbf{10.323,61 EUR} \\
\bottomrule
\end{tabularx}
\end{table}

Por tanto, el coste económico estimado de OmniLitter asciende a
\textbf{10.323,61 EUR} incluyendo contingencia. Esta cifra debe interpretarse como
una valoración del proyecto, no como un gasto real efectuado. El desembolso directo
fue muy reducido, pero el coste imputado refleja mejor la magnitud del trabajo
realizado y permite comparar el proyecto con un desarrollo software equivalente en
un entorno profesional.
